% Transcription — fonds Alexandre Grothendieck, shelfmark 12, batch 6 (pages 101-120).
% Established from the facsimile archives/batches/12/batch-06.pdf (a mirror of
% https://grothendieck.umontpellier.fr/12.pdf), per the skill
% transcribe-grothendieck. The batch file carries no cover sheet: PDF page N
% is page 100+N, as the Montpellier footer printed on each PDF page confirms,
% and the archivists' pencil numbers bottom-left agree with that position
% wherever they are legible (101, 112-120 checked).
%
% Pass: Opus 5.5 (claude-opus-5-5), 2026-09-23 — first pass, unchecked against the pages by a human.
%
% What the batch is.
%   101      An isolated typed leaf in English (cycle classes: functoriality
%            squares for f^*, f_*; correspondences via Künneth and Poincaré
%            duality; "algebraic" homomorphisms). Formulas and diagrams inked
%            by hand into the typed blanks; it breaks off before the
%            composition it announces. Not continued on either side.
%   102, 109 Blank. Absent.
%   103      A loose sheet: weight/level tables for H^1 and H^2 (proper vs
%            smooth case), semi-simplicity of the pure pieces, a remark on
%            H^2_! and mixing of weights "in the wrong direction".
%   104-108  "Poids de la cohomologie ℓ-adique d'une variété singulière ou/et
%            non propre" (his boxed title, p. 104). A) constant coefficients:
%            (1) X smooth quasi-projective, X = X' - Y with Y a normal-crossings
%            divisor; Proposition 1 (weights of H^i_! in [0, i], levels) and
%            Corollaire 1 by duality for H^i (weights in [i, 2n]); (2) smooth
%            non-quasi-projective / non-separated; (3), (4) non-smooth, by
%            dévissage and cohomological descent. B) twisted coefficients
%            coming from R^j f_*Q_ℓ of a projective smooth f, degeneration by
%            Deligne; C) effective coefficients of weight and level j.
%   110-115  "Faisceaux ℓ-adiques admissibles et motifs virtuels" (his title,
%            p. 110), numbered 1., 2.1-2.3, 3.1-3.3, 4.1-4.2, 5.1-5.4, 6.:
%            "spécial de poids i" sheaves (subquotients of R^i f_*Q_ℓ, f proj.
%            smooth), Lemma 2.2 (sums, tensor, pull-back, dual twist,
%            Q_ℓ(-i)); "effectivement admissible à poids dans P" (3.1-3.3),
%            with levels (4.1-4.2); Théorème 5.1 (R^i f_! of a strictly
%            special sheaf is eff. admissible, weights ≤ i+j, levels table),
%            its proof by reduction to A) (1); Cor. 5.2 under the
%            semi-simplicity part of Tate's conjectures; Cor. 5.3 without it;
%            Cor. 5.4 over a field for R^i f_*; §6 opens on R^i f_* for an
%            immersion. P. 115 is the calculation sheet for that immersion
%            case (normal-crossings compactification, spectral sequence).
%   116      A typed administrative note to him dated 27 March 1968. Absent.
%   117      Weight/level tables for the E_2^{p,q} terms, cases i ≤ n-1 and
%            i ≥ n; (p, q) grids drawn, not redrawn.
%   118      Printed cover of Springer LNM 41, "Local Cohomology" (1967), with
%            a printer's stamp "20.3.68 1. Korrektur". Absent.
%   119      Calculation sheet on the back of p. 120: the refined tables, the
%            boxed bounds weight ≤ 2 Inf(i, n), level ≤ Inf(i, n-1), then
%            + j for F special of weight j, and the relative case with
%            n = d_y(f).
%   120      An isolated typed leaf in French, starting and ending mid-sentence,
%            presenting Deligne's work (item 5: λ-structure on K^•(X) via
%            perfect complexes and Dold-Puppe; the SGA 6 seminar of 1966/67
%            is named). Hand additions « resté », « sur X », « -Puppe ».
%   The two skipped versos (116, 118) both carry March 1968 dates; they are
%   recorded here only, not in the body, since the skill skips such pages.
%
% Notation fixed for the batch. \operatorname{Inf} for his « Inf » (minimum),
% H^{i}_{!} for compactly supported cohomology, \mathbb{Q}_{\ell,X} for the
% constant sheaf, E_{2}^{pq} for spectral-sequence terms, \check{F} for the
% dual. His « eff^t » is eff\supplied{t} (effectivement), « str^t » str\supplied{t}
% (strictement); his « ½simplicité » is kept as \tfrac{1}{2} and glossed once.
% Prose ellipses are « … »; \ldots only inside math.
%
% Legibility. 104-107 and 110-111 are a fast draft but carry; 106, 108, 112,
% 113 are heavily struck over and reworked in the interlinear space — the
% formulas and the final state are given, much connecting prose is \ill{}.
% 115, 117, 119 are calculation sheets, transcribed as tables and formulas.
\documentclass[11pt,a4paper]{article}
% Le préambule commun à toutes les transcriptions du fonds Grothendieck.
%
% Il tient en une idée : **l'incertitude se compose, elle ne se résout pas.**
% Une transcription qui lisse un mot douteux a détruit la seule chose pour
% laquelle on la faisait — un lecteur doit pouvoir savoir, à chaque endroit,
% ce qui était sur le feuillet et ce qui a été ajouté. Les six macros qui
% suivent sont donc l'appareil critique, et non de la décoration.
%
% Les mêmes commandes sont reconnues par `scripts/render.mjs`, qui produit la
% vue de lecture du site. Une macro ajoutée ici doit l'être aussi là-bas,
% faute de quoi le rendu échouera bruyamment — ce qui est voulu : un rendu
% silencieusement faux est pire qu'un rendu absent.

\ProvidesPackage{grothendieck}[2026/08/08 Transcriptions du fonds Grothendieck]

\usepackage{amsmath,amssymb,amsthm}
\usepackage{mathrsfs}
\usepackage{stmaryrd}
% Les diagrammes commutatifs sont partout dans ce fonds : c'est la langue
% même de la géométrie algébrique de Grothendieck, et une transcription qui
% les rendrait en prose ne transcrirait rien.
\usepackage{tikz-cd}
% Français par défaut — c'est la langue des feuillets — mais l'anglais est
% chargé aussi : la traduction et le résumé sont des documents anglais, et la
% typographie française (espaces avant les deux-points) y serait une faute.
% Ces documents basculent par \selectlanguage{english} en tête de corps.
\usepackage[english,french]{babel}
\usepackage{fontspec}
\usepackage{geometry}
\geometry{a4paper,margin=25mm,marginparwidth=28mm}
\usepackage{marginnote}
\usepackage{xcolor}
% ulem, not soul. `soul`'s \st and \ul are built on a character-by-character
% scan that XeLaTeX's font machinery breaks: the document fails with an
% undefined control sequence rather than degrading. ulem does the same two jobs
% and is Unicode-engine safe. `normalem` stops it hijacking \emph, which the
% transcriptions use for Grothendieck's own emphasis.
\usepackage[normalem]{ulem}
\usepackage{hyperref}
\hypersetup{colorlinks=true,linkcolor=black,urlcolor=black,pdfborder={0 0 0}}

% --- Métadonnées du lot -----------------------------------------------------
% Reprises telles quelles par le script de rendu, qui les affiche en tête de
% la vue de lecture. Elles fixent ce que le fichier prétend couvrir : sans
% elles, une transcription flotte sans point d'attache dans un fonds de seize
% mille feuillets.

\newcommand{\folder}[1]{\gdef\grothFolder{#1}}
\newcommand{\batch}[1]{\gdef\grothBatch{#1}}
\newcommand{\leaves}[2]{\gdef\grothFirst{#1}\gdef\grothLast{#2}}
\newcommand{\foldertitle}[1]{\gdef\grothFoldertitle{#1}}
\gdef\grothFolder{?}\gdef\grothBatch{?}\gdef\grothFirst{?}\gdef\grothLast{?}\gdef\grothFoldertitle{}

% --- Le résumé --------------------------------------------------------------
%
% Il ouvre la lecture modernisée, sous le titre « Résumé », et il est composé
% en petit. Ce n'est pas un ornement : c'est le seul passage du document qui
% ne soit pas des mathématiques, et un lecteur doit pouvoir le reconnaître
% avant de l'avoir lu — puis le sauter s'il connaît déjà le sujet, ou s'y
% arrêter s'il ne le connaît pas. Le filet qui le suit marque la reprise du
% corps.

\newenvironment{resume}
  {\small\setlength{\parskip}{0.45em}}
  {\par\medskip\hrule\medskip}

% --- Mots-clés ---------------------------------------------------------------
%
% En anglais, en clôture du résumé de la lecture modernisée : le vocabulaire
% moderne sous lequel chercher ce que les feuillets construisent. Le manifeste
% du site (`npm run manifest`) les extrait pour étiqueter la cote entière —
% cette ligne est la source unique des tags, il n'y a pas de fichier à côté.
\newcommand{\keywords}[1]{\par\smallskip\noindent
  {\footnotesize\textsc{Keywords} --- \textit{#1}}\par}

% --- L'appareil critique ----------------------------------------------------

% Le numéro de feuillet, en marge. C'est l'ancre qui permet de régler un
% désaccord sur une formule en regardant *un* feuillet plutôt qu'une chemise
% de six cents. Le site s'en sert aussi pour tourner les pages du fac-similé
% quand on fait défiler la transcription.
\newcommand{\leaf}[1]{%
  \marginnote{\footnotesize\color{gray!70}\textsf{#1}}%
  \label{leaf:#1}\ignorespaces}

% Illisible : on ne devine pas. Un mot inventé qui se lit comme les autres est
% le pire résultat possible.
\newcommand{\ill}{\textcolor{red!60!black}{[\,\ldots\,]}}

% Lecture douteuse : le mot est donné, et signalé comme incertain.
\newcommand{\uncertain}[1]{\uline{#1}}

% Ajout de l'éditeur — une lettre manquante, un mot sous-entendu.
\newcommand{\add}[1]{\textcolor{blue!50!black}{[#1]}}

% Biffé par Grothendieck lui-même. Ce qu'il a rayé fait partie du document :
% une rature dit souvent où la pensée a changé de direction.
\newcommand{\struck}[1]{\sout{#1}}

% Note du transcripteur, sur l'état matériel du feuillet.
\newcommand{\note}[1]{\par\smallskip{\small\color{gray!60!black}\textit{[#1]}}\par\smallskip}

% Note marginale de Grothendieck, distincte du corps du texte.
\newcommand{\marginal}[1]{\par\smallskip\noindent\hspace{1em}%
  {\small\color{gray!60!black}#1}\par\smallskip}

% --- Titre ------------------------------------------------------------------

\AtBeginDocument{%
  \begin{center}
    {\large\bfseries Fonds Alexandre Grothendieck --- cote n\textsuperscript{o}~\grothFolder}\\[2pt]
    {\itshape\grothFoldertitle}\\[4pt]
    {\small feuillets \grothFirst--\grothLast\ (lot \grothBatch)}\\[2pt]
    {\footnotesize\color{gray!60!black}Transcription établie d'après le fac-similé numérisé
     par l'Université de Montpellier.\\
     Les lectures incertaines sont soulignées, les illisibles notées [\,\ldots\,],
     les ajouts entre crochets.}
  \end{center}
  \vspace{1em}}

\endinput


\folder{12}
\batch{6}
\pages{101}{120}
\dating{1964-[vers 1971]}
\watermark{Édition de démonstration}
\foldertitle{Généralités (sous-groupes de Galois motiviques) (1964 ?) : notes manuscrites (s.d.), tapuscrit (s.d.)}

\begin{document}

\page{101}
\note{feuillet dactylographié, en anglais, isolé : ni le feuillet qui le
précède dans le lot 5 ni celui qui le suit ici n'en continuent le texte, et
il s'arrête au milieu d'une phrase. Les formules et les deux diagrammes sont
tracés à la main dans les blancs laissés par la machine, de même que les
mentions « Kunneth formula » et « Poincaré duality » ; nous ne savons pas
de quelle main}

The association of a cohomology class with a cycle has all the functorial and
compatibility properties one could hope for. In particular, if
$f : X \to Z$ is a morphism of schemes over $k$, then there are commutative
diagrams

\begin{tikzcd}
H^{2i}_{\ell}(Z) \arrow[r, "f^{*}"] & H^{2i}_{\ell}(X) \\
C^{i}(Z) \arrow[u, "\rho_{Z}"] \arrow[r, "f^{*}"] & C^{i}(X) \arrow[u, "\rho_{X}"']
\end{tikzcd}

\begin{tikzcd}
H^{2i}_{\ell}(X) \arrow[r, "f_{*}"] & H^{2i-2r}_{\ell}(Z) \\
C^{i}(X) \arrow[u, "\rho_{X}"] \arrow[r, "f_{*}"] & C^{i-r}(Z) \arrow[u, "\rho_{Z}"']
\end{tikzcd}

where $r = \dim X - \dim Z$.

\emph{Correspondences.} Let $X$ and $Z$ be projective, smooth, and connected
schemes over $k$. Then there are canonical isomorphisms
\[
\begin{aligned}
H^{*}_{\ell}(X \times Z) &\simeq H^{*}_{\ell}(X) \otimes_{\mathbb{Q}_{\ell}} H^{*}_{\ell}(Z)
&& \text{(Kunneth formula)} \\
&\simeq \check{H}^{*}_{\ell}(X) \otimes_{\mathbb{Q}_{\ell}} H^{*}_{\ell}(Z)
&& \text{(Poincaré duality)} \\
&\simeq \operatorname{Hom}_{\mathbb{Q}_{\ell}}\bigl(H^{*}_{\ell}(X), H^{*}_{\ell}(Z)\bigr) .
\end{aligned}
\]
Thus to an element $\alpha$ of $H^{*}_{\ell}(X \times Z)$ we can canonically
associate a homomorphism $\tilde{\alpha} : H^{*}_{\ell}(X) \to H^{*}_{\ell}(Z)$,
and to a homomorphism $\beta : H^{*}_{\ell}(X) \to H^{*}_{\ell}(Z)$ we can
canonically associate an element $\hat{\beta} \in H^{*}_{\ell}(X \times Z)$.
In particular, we say that $\beta : H^{*}_{\ell}(X) \to H^{*}_{\ell}(Z)$ is
\emph{algebraic} if the associated cohomology class $\hat{\beta}$ is
\emph{algebraic}, i.e.\ belongs to $C^{*}_{\ell}(X \times Z)$.

It is sometimes convenient to note that, given
$\alpha \in H^{*}_{\ell}(X \times Z)$, the homomorphism
$\tilde{\alpha} : H^{*}_{\ell}(X) \to H^{*}_{\ell}(Z)$ is the composition of
the homomorphisms
\note{la composition annoncée n'est pas inscrite : le bas du feuillet est laissé en blanc}

\page{103}
\note{feuillet isolé, à l'encre bleue : deux tableaux de poids et de niveaux,
puis deux remarques ; la disposition en colonnes est celle de la page}

$H^{1}$ \quad \struck{\ill{}} \emph{cas propre, ou du $H^{1}_{!}$} (au-dessus
des colonnes de poids 0 et 1) ; \emph{cas lisse} (sous les colonnes de poids 1
et 2) :
\[
\begin{array}{c|ccc}
\text{poids} & 0 & 1 & 2 \\
\hline
\text{niveau} & 0 & 1 & 0
\end{array}
\]
morceaux purs sont semi-simples

$H^{2}_{\ill{}}$ \quad \emph{cas propre} (au-dessus des poids 0 à 2) ;
\emph{cas lisse} (sous les poids 2 à 4) :
\[
\begin{array}{c|ccccc}
\text{poids} & 0 & 1 & 2 & 3 & 4 \\
\hline
\text{niveau} & 0 & 1 & 2 & 1 & 0
\end{array}
\]
\note{l'indice de ce $H^{2}$ est une petite boucle que nous ne savons pas
lire ; ce n'est pas le « ! » de la ligne suivante}

morceaux purs de poids $\neq 2$ sont semi-simples. Le morceau pur de poids 2
est une extension par un motif de Tate $\mathbb{Q}(-1)^{r}$.

$H^{2}_{!}$ \quad il semble qu'il puisse y avoir mélange de poids dans le
mauvais sens, à cause d'une extension d'un motif effectif de poids 0 par un
autre de poids 1 ; en plus de cela, il y aurait une extension, peut-être non
triviale, d'un motif effectif de poids 0 \add{resp.} 1\note{« resp. 1 » est
ajouté au-dessus de la ligne} par un autre.

\section{Poids de la cohomologie $\ell$-adique d'une variété singulière ou/et non propre}
\note{titre encadré, en tête de la page 104, de sa main ; « $\ell$-adique » y
est ajouté au-dessus de la ligne, et une parenthèse, en haut à droite, porte
\uncertain{(sur un corps de type fini)}}

\page{104}
\subsection{A) Coefficients constants}

(1) Soit d'abord $X$ lisse et quasi-projective, donc par résolution forte (que
nous admettons), on aura
\[
X = X' - Y, \qquad Y = \bigcup_{i} Y_{i},
\]
$X'$ lisse et projective, $Y_{i}$ diviseurs lisses, irréd.\add{uctibles}, à
croisements normaux. \struck{Soit} On a donc
\[
H^{i-1}(X') \to H^{i-1}(Y) \to H^{i}_{!}(X) \to H^{i}(X') \to H^{i}(Y) \to
H^{i+1}_{!}(X) \to \cdots
\]
\note{le début de la suite, $H^{i-1}(X') \to H^{i-1}(Y)$, est récrit au-dessus
d'un premier terme biffé}
\note{en marge gauche, le diagramme}

\begin{tikzcd}
X \arrow[r, hook] \arrow[d, "f"'] & X' \arrow[dl, "f^{\prime}"] \\
\operatorname{Spec} k &
\end{tikzcd}

\struck{\ill{}} donc les poids dans $H^{i}_{!}(X)$ [qui se définissent
\uncertain{via les val.\ absolues des valeurs propres de Frobenius} dans le
cas \uncertain{d'un} corps fini, par réduction au cas d'un corps de type
fini…] sont contenus dans ceux de $H^{i}(X')$ [\uncertain{voisins} à $i$
par Weil] et ceux de $H^{i-1}(Y)$.

Or
\[
H^{*}(Y) \Longleftarrow E_{2}^{pq} = H^{p}\bigl(i_{0} \ldots i_{p} \mapsto
H^{q}(Y_{i_{0}} \cap \cdots \cap Y_{i_{p}})\bigr)
\]
\note{sous le terme, en interligne : « on peut se borner ici aux cochaînes
alternées » ; à droite, une accolade porte : de poids $q$ ; de niveau
$\leqslant q$ [intervient si $i (= p+q) \leqslant n-1$] ; de niveau
$\leqslant 2n-i-q$ [si $i = p+q \geqslant n$]. En marge gauche :
$(2n-i-1)-q$}

\struck{Donc} Or ici $Y_{i_{0}} \cap \cdots \cap Y_{i_{p}}$ est projectif et
lisse (de dimension $n-p-1$ si $X$ \struck{\ill{}} purement de dim.\ $n$).
Donc le poids du terme $E_{2}^{pq}$ est $q$, et on trouve que le poids de
$H^{j}(Y)$ est $\leqslant j$. Donc :

\textbf{Proposition 1} \struck{Lemme}\note{« Proposition 1 » est écrit
au-dessus d'un premier mot biffé} $H^{i}_{!}(X)$ est de poids compris entre 0
et $\leqslant i$. La partie de poids $i$ (resp.\ de poids $\leqslant i-1$) est
le quotient image de $H^{i}_{!}(X) \to H^{i}(X')$ [resp.\ l'image de
$H^{i-1}(Y) \to H^{i}_{!}(X)$]. De plus, si $i \geqslant n = \dim X$, alors le
\emph{niveau} de $H^{i}_{!}(X)$ est $\leqslant 2n-i$, et les deux niveaux
précédents sont de niveau $\leqslant 2n-i-1$ ; de plus, il n'y a des poids
qu'entre $2(i-n)$ [car il faut écrire $q \leqslant 2(n-p-1)$ !] et $i$, i.e.
\ill{} d'amplitude $2n-i$…

Pour un complément, notons que le niveau de $E_{2}^{pq}$ est $\leqslant$ celui
de $H^{q}(Y_{i_{0}} \cap \cdots \cap Y_{i_{p}})$ (de dim.\ $n-p-1$), donc
\[
\leqslant \struck{\ill{}}\; 2(n-p-1) - q \;=\; \struck{\operatorname{Inf}(q,}
(2n-i-2) - p \qquad (i = p+q),
\]
donc $\leqslant 2(n-1)-i$.
\note{au-dessus, une variante : « donc $= 2(n-1)-i-p$ »}

\marginal{NB Tous les niveaux de la filtration sont « effectifs »}
\marginal{En tous cas, niveau $\leqslant \operatorname{Inf}(i, n)$}

\page{105}
Donc \struck{remplaçant} pour $p+q = i-1$, on trouve un niveau
$\leqslant 2(n-1)-i+1 = 2n-i-1$.

\emph{Remarque.} La filtration de $H^{*}(Y)$ déduite de la suite spectrale
est aussi autre que la filtration par les poids. La filtration image de
$H^{i}_{!}$ est donc aussi la filtration par les poids croissants (de 0 à $i$,
\ill{}) donc est également indépendante de la \uncertain{compactification}
$X'$ choisie. Écrivons les termes de cette filtration, en utilisant le fait que
$H^{i-1}(X')$ est purement de poids $i-1$, et \struck{$E_{2}^{i-1,0}$} que la
suite spectrale est dégénérée pour des raisons de poids :
\[
\begin{array}{c|c|c|c|c|c|c}
 & E_{2}^{i-1,0} & E_{2}^{i-2,1} & \cdots & E_{2}^{1,i-2} &
 E_{2}^{0,i-1}/\operatorname{Im} H^{i-1}(X') &
 \operatorname{Ker}\bigl(H^{i}(X') \to E_{2}^{0,i}\bigr) \\
\hline
\text{poids} & 0 & 1 & & i-2 & i-1 & i
\end{array}
\]
\note{sous le tableau, une troisième ligne est amorcée puis laissée vide, son
intitulé biffé}

\struck{Cas $i \geqslant n$} Si $i \leqslant n$, les « niveaux » sont
\add{seulement} $\leqslant$ poids, il n'y a pas lieu à priori d'attendre plus.
L'on trouve d'autre part, pour $i \geqslant n$, un tableau
\uncertain{de gr.\ filtré $\geqslant n$} :
\[
\begin{array}{c|c|c|c|c|c|c}
 & 0 \cdots 0 & \scriptstyle E_{2}^{2n-i-1,2(i-n)} & \scriptstyle E_{2}^{2n-i-2,2(i-n)+1} & \cdots
 & \scriptstyle E_{2}^{0,i-1}/\operatorname{Im} H^{i-1}(X') &
 \scriptstyle \operatorname{Ker}(H^{i}(X') \to E_{2}^{0,i}) \\
\hline
\text{poids} & \times & 2(i-n) & 2(i-n)+1 & \cdots & i-1 & i \\
\hline
\text{niveau} & \times & 0 & 1 & \cdots & 2n-i-1 & 2n-i
\end{array}
\]
\note{les premières colonnes, nulles, sont barrées d'une croix sur les deux
lignes ; l'exposant de la première colonne non nulle est récrit sur un premier
exposant biffé}

\textbf{Corollaire 1} (\uncertain{par dualité}) $H^{i}(X)$ est de poids
compris entre $i$ et $2n$. Le morceau \struck{\ill{}} de poids $i$ est
exactement l'image de $H^{i}(X')$. Si $i \geqslant n$, les niveaux des morceaux
de poids $i, i+1, \ldots, 2n$ sont resp.\ $\leqslant 2n-i, 2n-i-1, \ldots, 0$.
Si $i \leqslant n$, il n'y a de morceaux que des poids $i, i+1, \ldots, 2i$
\struck{$i, i+1, \ldots, 2n = 2(n-i) = \ldots$}, qui sont de niveaux resp.
$\leqslant i, \uncertain{i-1}, \ldots, 0$.
\marginal{NB Les niveaux ici \emph{dé}croissants}

\marginal{Tous les morceaux de la filtration sont « effectifs ». En tous cas,
niveau $\leqslant \operatorname{Inf}(i, n)$}

\page{106}
(2) \emph{Cas \struck{\ill{}} lisse non quasi-projectif.}

Si $X$ est séparé, utilisant une compactification de Nagata, et résolution des
singularités [comme ci-dessus], on a les mêmes résultats que ci-dessus, en
notant que la résolution des singularités permet de ramener les conj.\ de Weil
pour des var.\ propres et lisses au cas projectif et lisse.

Si $X$ non séparé, $H^{i}_{!}(X)$ n'est plus défini. Par contre $H^{i}(X)$ est
défini. On a, avec un recouvrement affine,
\[
H^{*}(X) \Longleftarrow E_{2}^{pq} = H^{p}\bigl(\sigma \mapsto
H^{q}(U_{\sigma})\bigr) .
\]
Ici $E_{2}^{pq}$ est effectif de poids $\geqslant q$, mais
\struck{non dégénéré} \struck{\ill{}} on peut en conclure que l'effectivité de
$H^{i}(X)$, et le fait que ses poids sont compris entre 0 et $2n$ si
$i \geqslant n$, et entre 0 et $2i$ si $i \leqslant n$ ; en tous cas, en
\ill{}, on trouve que le niveau \struck{est $\leqslant 2n-i$ (il reste à
regarder pour $i \geqslant n$)} \uncertain{des poids $\leqslant 2i$} est
$\leqslant i$ \struck{\ill{}}.
\note{passage très raturé : deux lignes biffées d'un long trait ; nous ne
donnons que ce qui reste lisible}
\marginal{Question : le niveau $\leqslant 2n-i$ ? \struck{\ill{}} (si
$n = \dim X$)}

(3) \emph{Cas non lisse} : $H^{i}_{!}(X)$, ($X$ séparé). Je dis qu'il est encore
\emph{effectif} et de poids $\leqslant i$, de niveau
$\leqslant \operatorname{Inf}(i, 2n-i)$ \uncertain{le complément ci-dessus}.
\note{en interligne : « vrai quand il est lisse dense »,
\uncertain{lecture douteuse}}
Soit $U$ un ouvert lisse dense, $Y = X - U$. On a
\[
\to H^{i}_{!}(U) \to H^{i}_{!}(X) \to H^{i}_{!}(Y) \to
\]
et $H^{i}_{!}(U)$ est effectif de poids $\leqslant i$, \add{de niveau
$\leqslant i$}\note{ajouté au-dessus de la ligne}, et $H^{i}_{!}(Y)$ aussi par
hypothèse de récurrence sur la dimension, donc $H^{i}_{!}(X)$ est effectif de
poids $\leqslant i$, de niveau $\leqslant i$. De même $H^{i}_{!}(U)$ est de
niveau $\leqslant 2n-i$, et $H^{i}_{!}(Y)$ de niveau
$\leqslant 2(n-1)-i \leqslant 2n-i$, cqfd.

(4) \emph{Cas non lisse} : $H^{i}(X)$ \struck{($X$ séparé)}. Je dis qu'il est
effectif, à poids $\leqslant 2i$ \add{$\leqslant i$ et $\leqslant n$}, à
niveau \struck{$\leqslant \operatorname{Inf}(i, \ldots\ \dim X)$} \ill{}.
\note{l'ajout « $\leqslant i$ et $\leqslant n$ » est écrit au-dessus de la
ligne, entre « poids » et « niveau », sans qu'on voie auquel des deux il se
rapporte}
Par la suite spectrale de Leray, on est ramené au cas où $X$ affine. Récurrence
sur $i$, c'est trivial si $i \leqslant 0$. On utilise une résolution
$X' \to X$, et suite spectrale de descente
\[
H^{*}(X) \Longleftarrow E_{2}^{pq} = H^{p}\bigl(s \mapsto
H^{q}((X'/X)^{s})\bigr)
\]
\note{l'exposant $s$ porte sur le produit fibré itéré de $X'$ au-dessus de
$X$ : \uncertain{lecture}}
par hyp.\ de récurrence, si $q < i$, $E_{2}^{pq}$ est effectif
\marginal{poids $\leqslant 2\operatorname{Inf}(i,n)$, niveau $\leqslant
\uncertain{\operatorname{Inf}(i,n)}$}

\page{107}
de \struck{niveau} poids $\leqslant 2q$, de niveau $\leqslant q$. Donc il en
est ainsi de $E_{\infty}^{pq}$ si $p+q = i$, $p \neq 0$, \add{et même ils
sont de poids $\leqslant 2(i-1) < 2i$, de niveau $\leqslant i-1 < i$}\note{en
interligne, au-dessus de la ligne} ; reste à examiner le terme
$E_{2}^{0,i} \subset H^{i}(X')$ ; comme $X'$ est lisse, il est donc effectif,
de poids $\leqslant 2i$, de niveau $\leqslant i$. On ajoute que, si $i > n$,
on a $H^{i}(X) = 0$ ($X$ affine !).

\struck{A fortiori $H^{i}(X)$ est de niveau $\leqslant n = \dim X$. Il suffit
\ill{} $E_{2}^{pq}$ \ill{} utilisant un recouvrement}
\note{trois lignes cancellées d'un long zigzag, lisibles en partie}

\subsection{B) Coefficients tordus sur $X$ lisse}
\note{dans le titre, deux ajouts entourés : « séparé » après « lisse »,
et « $F$ effectif de poids et niveau $j$ »}

(1) On suppose que les coefficients envisagés proviennent d'un motif lisse sur
$X$, et \add{par dévissage (possible loc.\ ?)}\note{ajouté en interligne} on se
ramène au cas où ce dernier est un facteur direct d'un
$R^{j}f_{*}(\mathbb{Q}_{\ell,Y})$, avec $f : Y \to X$ un morphisme projectif
et lisse. \add{On peut supposer même $f$ de dim.\ relative $j$.}\note{ajouté
en interligne, suivi d'une ligne que nous ne lisons pas} On a
\[
F = R^{j}f_{*}(\mathbb{Q}_{\ell,Y}) .
\]
Considérons la suite spectrale
\[
H^{*}(Y) \Longleftarrow E_{2}^{pq} = H^{p}\bigl(X, R^{q}f_{*}(Y)\bigr) .
\]
D'après \emph{Deligne}, cette suite spectrale dégénère. Donc
$H^{i}(X, R^{j}f_{*}(\mathbb{Q}_{\ell}(Y)))$ \add{et par suite $H^{i}(X,F)$}
apparaît comme un \emph{morceau} de $H^{i+j}(Y)$. On trouve que c'est un
\struck{module} (effectif) de poids compris entre $i+j$ et
$2(n+j) = 2n+2j$ (où $n = \dim X$) ; \struck{$2(n+j)$ et $2(i+j)$} les
niveaux sont resp.\ $\leqslant (2n-i+j), \ldots, 0$. Si $i \leqslant n$, il
n'y a que des poids \uncertain{qui ont les dimensions}
$i+j, i+j+1, \ldots, 2(i+j)$, de \struck{poids} niveaux respectifs
\struck{$\leqslant 0, \ldots, i+j$} $i+j, i+j-1, \ldots, 0$.
\note{« $2(i+j)$ » et « $2n+2j$ » sont entourés ; les trois bulles de la marge
s'y rattachent}
\marginal{peut-on remplacer $2n+2j$ par $2n+j$ ?}
\marginal{peut-on remplacer $2i+2j$ par $2i+j$ ? \struck{Idiot \ill{}} !}
\marginal{Les poids sont-ils $\geqslant j$ ?}

Résultats duals pour $H^{i}_{!}(F)$, il est effectif de poids $\leqslant i+j$,
et \struck{à niveaux} si $i \geqslant n$, il est à niveaux croissants de 0 à
\struck{$2n-i+j$} $2(n+j)-(i+j) = 2n-i+j$ pour des poids croissants de
\ill{}.

Si $X$ est \emph{propre}, on trouve que $H^{i}(X,F)$ est pur de poids $i+j$
(et semi-simple ici…) !

\page{108}
(2) Cas où on dispose seulement d'une représentation de $F$ sous forme
symbolique \emph{localement} sur $X$ pour la top.\ (Zariski). Dans ce cas,
\struck{on ne peut plus utiliser la filtration. $X$ séparé, on doit
\ill{} localiser sur $X$ en sous (ét), et utiliser la suite spectrale de Leray
d'un recouvrement. On trouve \ill{} effectivité, poids entre 0 et $2n+2j$}
\note{passage encadré à gauche, biffé de deux traits obliques et d'un trait
horizontal} on a intérêt à travailler \emph{d'abord} avec la cohomologie à
supports propres sur $X$, qui permet de découper $X$ en morceaux, et d'obtenir
le même résultat que dessus. On déduit l'analogue en cohomologie à supports
quelconques par dualité.

\struck{Si $F$ est symplectique,}

\emph{Remarque.} [En fait, on voit par ce raisonnement que, en ce qui concerne
la coh.\ à supports propres, les résultats sont obtenus en supposant seulement
que $F$ soit « joli par morceaux » sur $X$ \add{\uncertain{la lissité de $X$
par ailleurs}} ; si \uncertain{de plus} $X$ lisse, le fait \uncertain{que}
$F$ \uncertain{constant} est utilisé \struck{\ill{}} pour passer aux
\uncertain{résultats} précédents par dualité.

\subsection{C) Cas de coefficients effectifs de (poids $j$, niveau $j$) quelc., $X$ quelconque}

1) Cas de la cohomologie $H^{i}_{!}$ : déjà traité dans le numéro
précédent : résultat le même que dans (B) 1).

2) Cas de la cohomologie $H^{i}(X, F)$. Je dis qu'il est \emph{effectif}, de
poids \struck{entre 0 et $2i$}
$\leqslant \operatorname{Inf}(2i, 2n) = 2 \operatorname{Inf}(i, n)$, de niveau
$\leqslant \operatorname{Inf}(\uncertain{i}, n)$, où $n = \dim X$.
\struck{On procède par} ----
\marginal{?}

\section{Faisceaux $\ell$-adiques admissibles et motifs virtuels}
\note{titre en tête de la page 110, de sa main}

\page{110}
\struck{1. Soit $S$ un schéma, $\mathbb{Q}_{\ell}$}

1. On a fixé un nb premier $\ell$. Les schémas envisagés sont premiers à
$\ell$. (Le cas le plus intéressant est celui des schémas loc.\ de type fini
sur $\mathbb{Z}$.)\note{« de type fini sur » est récrit au-dessus de la
ligne ; un mot ajouté, \uncertain{constructibles}, s'y rattache} Les faisceaux
envisagés sont les $\mathbb{Q}_{\ell}$-faisceaux. On note $\mathbb{Q}_{\ell}(1)$
le faisceau de Tate $T_{\ell}(\mathbb{G}_{m}) \otimes_{\mathbb{Z}_{\ell}}
\mathbb{Q}_{\ell}$, et le twist par sa puissance tensorielle $n$-ième est
noté $F \mapsto F(n)$.
\marginal{Le bon point de vue \uncertain{commence} à se \uncertain{voir} !}

2.1 Un faisceau $F$ sur $S$ est dit \emph{spécial} \add{de poids $i$} s'il
est isomorphe \struck{à un facteur direct} \add{quotient tordu}
\add{sous-quotient tordu} d'un faisceau de la forme $R^{i}f_{*}(\mathbb{Q}_{\ell,X})$,
où $f : X \to S$ est projectif et lisse. [On entend \uncertain{que} si $i < 0$,
cela signifie $F = 0$ ; si $i = 0$, alors \struck{\ill{}} cela signifie :
$F$ localement constant.]
\note{deux ajouts entourés, dans l'interligne : « \uncertain{ou d'un
\ill{} ?} » et « sous-quotients et quotients tordus, et sommes directes » ; ce
dernier est rattaché par un trait au mot « équivalent à »,
\uncertain{lecture}}
\marginal{définition canulée : \uncertain{on ne saurait} se ramener à une
situation définie sur $\mathbb{Z}$}
\marginal{Un travail dans la catégorie des faisceaux $\ell$-adiques naturels
\uncertain{arithmétiques}…}
\note{sous ces deux bulles de marge, une troisième, longue, est cancellée de
traits verticaux ; nous ne la lisons pas}

\textbf{Lemme 2.2.}
\begin{enumerate}
\item[a)] Stabilité de « spécial de poids $i$ » par sommes directes.
\item[b)] Si $F$, $F'$ sont spéciaux de poids respectifs $i, j$, alors
$F \otimes F'$ est spécial de poids $i+j$.
\item[c)] Stabilité par image inverse.
\item[d)] Si $F$ spécial de poids $i$, $\check{F}(-i)$ l'est aussi.
\item[e)] $\mathbb{Q}_{\ell}(-i)$ est spécial de poids $2i$, pour $i \geqslant 0$.
\end{enumerate}
\note{ajout en interligne sous d), rattaché à la marge « d) \uncertain{on
déduit} Lefschetz »}

Démonstration.
\begin{enumerate}
\item[a)] Se prouve en introduisant $\coprod_{1 \leqslant i \leqslant n} X_{i}$,
\item[b)] Se prouve en considérant un $X_{1} \times_{S} X_{2}$, et utilisant
Künneth,
\item[c)] Par th.\ de changement de base.
\item[d)] Par th.\ de Lefschetz en cohomologie $\ell$-adique.
\item[e)] On prend $\mathbb{Q}_{\ell}(-i) = R^{2i}f_{*}(\mathbb{Q}_{\ell,
\mathbb{P}^{i}_{S}})$.
\end{enumerate}
\marginal{Cor.\ 2.3 Si $F$ spécial de poids $i$, et $n \geqslant i$, alors
$\check{F}(-n)$ est spécial de poids $2n-i$}

3.1 Un faisceau $F$ sur $S$ est dit \emph{effectivement admissible}
\struck{\ill{}} [\emph{admissible}] s'il existe pour tout ouvert $U$ de $S$
\struck{une filtration de $F|U$} \add{$F|U$ dont les quotients \ill{}} une
filtration …, où $Y \to X$ …
\struck{\ill{} et \ill{} $F|S_{\alpha}$ …}
\note{les lignes qui suivent la définition sont barrées de trois longs traits
et récrites en interligne ; nous ne donnons que les fragments lisibles}
\marginal{il faudrait permettre une extension finie radicielle loc.\ \ill{}
\uncertain{sur} $Y$ \uncertain{rendant} \ill{} spécial}

On dit que $F$ est \emph{à poids dans $P$}, où $P$ est une partie de
$\mathbb{Z}$, si … \uncertain{les} $F|S_{\alpha}$ sont
\uncertain{admissibles} … poids $\in P$. (Donc \struck{eff.\ admissible}
alors $F$ est de poids $\in P'$, pour tout $P' \supset P$.)
Si \struck{$P = \emptyset$} \struck{\ill{}} $P$ ne contient pas …
d'élément $\geqslant 0$, cela signifie $F = 0$.
\marginal{Attention, il est tautologique ! cela signifie que si $F$ est à
poids $\in P$ et à poids $\in P'$, il est à poids $\in P \cap P'$
\struck{\ill{}}}

\page{111}
\struck{d) Si $F$ spécial de poids}

\textbf{Proposition 3.3} \struck{3.2}
\begin{enumerate}
\item[a)] Stabilité de « eff.\ adm.\ de poids $\subset P$ » par
sous-faisceaux constructibles, quotients constructibles, \add{sommes
finies,} extensions.
\item[b)] Si $F, F'$ sont eff.\ adm.\ à poids $\subset P, P'$, alors
$F \otimes F'$ est admissible à poids
$\subset P + P' = \{ i+i' \mid i \in P, i' \in P' \}$.
\item[c)] Stabilité par chgt de base.
\item[d)] Si $F$ est \struck{spécial} eff\supplied{t} admissible \emph{constant
tordu} à poids $\subset P$, et si $i \in P \Rightarrow i \leqslant n$, alors
$\check{F}(\struck{n})$ est à\supplied{ eff\supplied{t}} admissible (constant tordu) à
poids $\subset 2n - P = \{ n-i \mid i \in P \}$.
\end{enumerate}
\note{sic : l'ensemble est écrit $\{ n-i \mid i \in P \}$, alors que
« $2n - P$ » appellerait $\{ 2n-i \}$ ; le twist de $\check{F}$ est
surchargé et nous ne le lisons pas avec sûreté}
\marginal{d) \uncertain{\ill{}} Lefschetz}

\emph{Dém.}\ a) b) c) sont immédiats, en utilisant les notations
correspondantes de 2.1. Notons d'ailleurs :

3.2. [\struck{Supposons $S$ \ill{}} \add{Supposons $S$ qu.\ cpt
\uncertain{sép}.}]\note{ajout entouré, au-dessus de la ligne} Pour que $F$
soit eff\supplied{t} adm.\ à poids $\subset P$, il faut et il suffit que l'on
puisse trouver une partition de $S$ en \struck{\ill{}} sous-schémas
$S_{i} \overset{\alpha_{i}}{\hookrightarrow} S$, $\alpha_{i}$ immersion de
prés.\ finie, avec \struck{\ill{}} $F|S_{i} = F_{i}$ admettant une filtration
dont les quotients successifs sont spéciaux de poids $\in P$.
\note{une grande accolade tracée à gauche, de 3.2 jusqu'en haut de la page,
se termine par une flèche vers « Proposition 3.3 » : l'énoncé 3.2 doit
précéder la Proposition}

Comme, pour $F$ constant tordu, la formation de $\check{F}$ commute au
changement de base, \add{3.2} d) résulte de 2.2.

4.1 Un faisceau $F$ est dit eff\supplied{t} admissible \emph{de niveaux}
$\leqslant \nu$ \struck{s'il est effectivement admissible, et} dans 3.1, on
peut prendre $G$ \struck{\ill{} de poids $j \in$} de la forme $G'(-k)$,
$k \geqslant 0$, où $G'$ est spécial de poids \struck{\ill{}} $\leqslant \nu$.
--- On dit que $F$ est eff\supplied{t} admissible « de poids $\subset P$, et de
niveau $\leqslant \nu$ » si, avec \struck{\ill{}} $G = G'(-k)$ dans 3.1, $G'$
est spécial de poids $j \leqslant \nu$, avec $j + 2k \in P$.
\emph{Attention}, à priori cela n'est pas la conjonction de « de poids
$\subset P$ » et « de niveau $\leqslant \nu$ ». \struck{\ill{}} On notera
que si \struck{\ill{}} $\nu$ est \struck{\ill{}} \struck{\ill{}} un entier tel
que $P \leqslant \nu$ (i.e.\ $j \in P \Rightarrow j \leqslant \nu$), alors
« eff\supplied{t} adm.\ de poids $\in P$ » $\Rightarrow$ « id.\ et à niveau
$\leqslant \nu$ ».
\note{en bas, avant « alors », un début de ligne biffé :
« \struck{4.2 a) b) c)} »}
\marginal{blague : après 3.2, on remarque qu'on peut faire : 3.2, soit 3.1
qu'on \uncertain{étend} par \uncertain{4.2}}

\page{112}
\textbf{Proposition 4.2}
\begin{enumerate}
\item[a)] Stabilité de « eff\supplied{t} adm.\ de poids $\subset P$ et niveau
$\leqslant \nu$ » par sous-faisceaux constr., quotients constructibles,
sommes finies, extensions.
\item[b)] $F_{i}$ « eff\supplied{t} adm.\ de poids $\subset P_{i}$, niveau
$\leqslant \nu_{i}$ » pour $i = 1, 2$ implique $F_{1} \otimes F_{2}$
« eff\supplied{t} adm.\ de poids $\subset P_{1} + P_{2}$, niveau
$\leqslant \nu_{1} + \nu_{2}$ ».
\item[c)] Stabilité par chgt de base.
\item[d)] Si $F$ \add{eff\supplied{t}} admissible \emph{constant tordu}, « à poids
\struck{$k$} \add{$\nu + 2k$} et niveau $\leqslant \nu$ » (où $k \geqslant 0$,
$\nu \geqslant 0$), \struck{et si} alors pour tout $n \geqslant \nu + k$
\add{(soit $n = \nu + k + h$, $h \geqslant 0$)}, $\check{F}(n)$ est « eff\supplied{t}
admissible constant tordu, à poids $2n - \nu - 2k = \nu + 2h$, niveau
$\leqslant \nu$ ».
\end{enumerate}
\note{le twist de $\check{F}$ est surchargé : lecture $\check{F}(n)$
\uncertain{probable}}

\emph{Remarque} : introduire plutôt : « à morceaux de poids $i_{\alpha}$,
niveau $\nu_{\alpha}$ », $\alpha \in \ldots$ ».

\textbf{Théorème 5.1.} Soit $f : X \to S$ un morphisme séparé de
\uncertain{type} \add{prés.\ finie}, et $F$ \struck{\ill{} de poids \ill{}}
\emph{strictement spécial} de poids $j$ \struck{(i.e.\ déf.\ 2.1, avec
\ill{})} (i.e.\ déf.\ 1 \add{\uncertain{d'un}} facteur \emph{direct} d'un
$R^{j}g_{*}(\mathbb{Q}_{\ell,Y})$, avec $g : Y \to X$ proj.\ lisse). Alors
\struck{\ill{}} sur $X$ \struck{\ill{}} \uncertain{s'appliquant}, avec $F$
…, $R^{i}f_{!}(F)$ est eff\supplied{t} admissible \struck{\ill{}} à poids
$\leqslant i + j$ \struck{\ill{}}, et à niveaux $\leqslant
\operatorname{Inf}(i+j, 2n-i+j)$ …
\struck{Si $f$ \ill{} lisse et $X$ \ill{} \ill{} poids \ill{} niveau
\ill{}} Si $n = \operatorname{Sup}_{s \in S} \dim X_{s}$, et si
$i \geqslant n$, les poids $\ldots$ possibles sont et niveaux correspondants
\[
\begin{array}{c|c|c|c|c}
\text{poids} & 2(i-n) & 2(i-n)+1 & \cdots & (i+j) \\
\hline
\text{niveau} & \leqslant 0 & \leqslant 1 & \cdots & \leqslant 2n-i+j
\end{array}
\]
\note{énoncé très raturé : deux lignes barrées, où figurait un premier
tableau (« poids $i+j$, $2i$ … ; niveau … ») biffé en zigzag ;
nous donnons l'état final, et le tableau est celui qu'il a gardé. La
première colonne, $2(i-n)$, ne dépend pas de $j$ : sic}
\marginal{Utiliser Lefschetz \uncertain{à la} Deligne}
\marginal{+ résolution des singularités pour les \uncertain{variétés sur}
les corps résiduels de $S$.}
\marginal{+ $\tfrac{1}{2}$ simplicité arithmétique ?}
\note{le $\frac{1}{2}$ de la dernière bulle : sa manière d'écrire
« semi- », soit « semi-simplicité
arithmétique ? »}

\emph{Démonstration.} Cas \add{(a)} où $F = R^{j}g_{*}(\mathbb{Q}_{\ell,Y})$,
$g : Y \to X$ proj.\ lisse. Utilisant le th.\ de changement de base propre,
\struck{\ill{}} et dévissage via suite exacte de cohomologie à supports
propres, on se ramène au cas où \struck{$\exists\, S' \to S$} $S$ est affine,
noeth., intègre, $\exists\, S' \to S$ morphisme fini \uncertain{radiciel}
\add{$X$ affine,} tel que $X_{1} = X \times_{S} S'$ soit \emph{lisse} sur
$S'$, et se « compactifie » en un $\widehat{X}_{S'}$ \struck{tel que $X'$}
projectif et lisse, tel que $\widehat{X}_{S'} - X_{S'}$ soit un diviseur
\add{relatif}, à croisements normaux, à composantes irréductibles
\uncertain{lisses} sur $S'$ ; enfin $F$ \struck{\ill{}} est \emph{spécial}
de poids $j$. [\struck{Quitte} : quitte à changer de notation, on
\uncertain{peut} … $S = S'$, et d'où $X \subset X' = \widehat{X}$.]
On \uncertain{peut supposer} $F = \mathbb{Q}_{\ell,X}$, on est ramené à
\uncertain{la} situation traitée dans les notes préliminaires (partie
(A, 1)). [\emph{NB} On n'utilise pas la dégénérescence de la suite
spectrale qui y intervient.]
\marginal{à expliciter au point \uncertain{de vue} \uncertain{construction}
servira dans 3.3.\ a) et 4.2 a)}
\note{la « partie (A, 1) » des « notes préliminaires » est la section
A) (1) des pages 104–105}

\page{113}
\struck{B) Cas général. On se réduit au}

Si $F$ est quelconque, on écrit que $F$ est un sous-…
\struck{facteur direct d'un} \add{ct.\ tordu} \struck{quotient}
\struck{d'un $R^{d}f_{*}(\mathbb{Q}_{\ell,Y})$}, avec $f : Y \to X$
morphisme projectif. [Utilisant Lefschetz \struck{et 2.1.\ d), …
$j \geqslant d$ … relative} \struck{\ill{} … la
\uncertain{fibre} de $f$} \struck{fortement} si $j < d =$ dim.\ relative de
$Y/X$, et remplaçant $Y$ par $Y \times_{X} \mathbb{P}^{j-d}_{X}$ si
$j > d$, on peut supposer $j = d$.]
\note{passage très raturé : trois lignes biffées, reprises dans
l'interligne ; nous donnons l'état final}
\marginal{\struck{\ill{} Lefschetz \ill{}} \uncertain{Voir si cette
réduction est vraiment utile}…}

La suite spectrale
\[
R^{*}_{!}(fg)(\mathbb{Q}_{\ell}) \Longleftarrow E_{2}^{pq} =
R^{p}_{!}f\bigl(R^{q}_{!}g(\mathbb{Q}_{\ell,X})\bigr)
\]
dégénère \struck{par} \emph{Deligne}, donc
$R^{i}_{!}f(F) = E_{2}^{ij}$ est $\simeq$ quotient \add{ct.\ tordu} d'un
sous-faisceau ct.\ tordu d'un $R^{i+j}_{!}(fg)(\mathbb{Q}_{\ell,Y})$.
\note{la suite spectrale est notée ici avec $f$ et $g$ ; le terme de gauche
porte $(fg)$ et $\mathbb{Q}_{\ell}$ sans indice lisible, et
$R^{q}_{!}g(\mathbb{Q}_{\ell,X})$ porte $X$ où l'on attendrait $Y$ : tel
quel}
La conclusion en résulte, en utilisant le cas précédent pour $Y \to S$, de
dim.\ relative $\leqslant n + j$…
\marginal{utiliser Lefschetz fort \uncertain{\ill{}} Deligne}
\marginal{[N.B.\ Voir si Deligne ne peut pas remplacer une factorisation
\uncertain{d'Artin} sur $X \to S$ \ill{} \uncertain{variant} …]}

\textbf{Corollaire 5.2.} Si la partie « \emph{semi-simplicité} » des
conjectures de Tate est vérifiée, la conclusion de 5.1 est vraie
\add{même si} on suppose seulement que $F$ est eff\supplied{t} admissible de poids
$j$. Par suite, si on suppose $F$ eff\supplied{t} admissible de poids
$\leqslant j$, alors on a un énoncé : $R^{i}_{!}f(F)$ est eff\supplied{t}
admissible, de poids \struck{\ill{}} $\leqslant i+j$, \add{niveau
$\leqslant \operatorname{Inf}(i+j, 2n-i+j)$} ; si $i \geqslant n$, il a des
\struck{\ill{}} poids \uncertain{entre} $2(i-n)$ et $(i+j)$, correspondant
aux niveaux $0, 1, \ldots, 2n-i+j$.
\marginal{semi-simplicité utilisée !}
\marginal{alors par semi-simplicité : $F$ eff\supplied{t} $=$ str.\ spécial}

\textbf{Cor.\ 5.3} \add{(Sans conj.\ de Tate)} On se ramène au cas
\add{par dévissage et déf.}\ où $F$ est \emph{spécial} de poids $j$, et par
\uncertain{th.\ de chgt de base propre} au cas où $S$ de type fini sur Spec
$\mathbb{Z}$, et \struck{\ill{}} si $F$ est eff\supplied{t} admissible, alors
$R^{i}_{!}f(F)$ est eff\supplied{t} admissible.
\note{l'énoncé de 5.3 et le début de sa réduction sont mêlés sur la page,
l'ajout « (Sans conj.\ de Tate) » et la réduction étant écrits au-dessus ;
nous suivons l'ordre des lignes}

\page{114}
\textbf{Cor.\ 5.4.} Sur un corps $k$, si $X$ est lisse, séparé, de t.f., et si
$F$ est \struck{\ill{}} str\supplied{t} spécial de poids $j$ (resp.\ eff\supplied{t}
adm.\ de poids $j$, si l'on admet conjecture de $\tfrac{1}{2}$simplicité),
alors $R^{i}f_{*}(F)$ est \struck{\ill{}} eff\supplied{t} admissible de poids entre
$i+j$ et $2n+2j$, et les niveaux correspondants étant resp.
$\leqslant (2n-i+j), 2n-i+j-1, \ldots, 0$. Si … $i \leqslant n$, on a
\struck{les niveaux} \emph{seulement} des poids $i+j, i+j+1, \ldots, 2(i+j)$,
\struck{dans les \ill{}} de niveaux correspondants resp.
$i+j, i+j-1, \ldots, 0$.
\note{« $\tfrac{1}{2}$simplicité » : semi-simplicité. Les niveaux
$2n-i+j$, $2n-i+j-1$ sont récrits sur des premiers termes surchargés}
\marginal{résolution, Lefschetz fort, et \uncertain{au besoin}
semi-simplicité}

On peut, sans procéder par dualité à partir de 5.1 resp.\ 5.2, \add{dans le
cas $F$ quelc.,} tout reprendre \add{(la dém.\ de 5.1 resp.\ 5.2,
\uncertain{via} la suite spectrale de Leray)}, puis se ramener
\struck{en utilisant la suite exacte de cohomologie relative} au cas où
$F = \mathbb{Q}_{\ell}$. Dans ce cas, on peut raisonner comme au cas
précédent par \add{compactification et} résolution des singularités, pour
écrire $X = \widehat{X} - D$, $\widehat{X}$ \struck{projectif et lisse} propre
et lisse et $D$ \add{diviseur} à croisements normaux, et on écrit suites
exactes de cohomologie \uncertain{relative} (supports quelc.)\ pour
$X, \widehat{X}$ et $D$, et la suite spectrale de localisation pour
$H^{*}_{D}(\widehat{X})$.
\marginal{Le faire à titre de vérification de cohérence du yoga !}

6. Soit $f : X \to S$ un morphisme de t.f., $F$ un faisceau
\struck{str\supplied{t} spécial} \add{effectivement admissible} sur $X$, de poids
$j$. On veut des renseignements sur \add{les} $R^{i}f_{*}(F)$, en
particulier on veut qu'ils soient effectivement admissibles, de poids que
nous préciserons. Supposons d'abord $F = \mathbb{Q}_{\ell,X}$.
\struck{Comme} par le th.\ de finitude, on regarde d'abord le cas d'une
immersion
\note{la page s'arrête sur ce mot}

\page{115}
\note{feuillet de calculs, sans phrases suivies, et qui semble continuer le
cas de l'immersion par lequel s'arrête la page 114 ; nous en donnons les
formules dans l'ordre de la page. Le haut porte quelques symboles
d'essai :}

$X \subset Y$ \qquad $X_{0}, X_{1}, \ldots$ \qquad \struck{$X \hookrightarrow Y$}
\qquad $\mathbb{Q}_{\ell}$ \qquad $X$ \emph{régulier}
\[
X_{0} \overset{i}{\hookrightarrow} X \overset{f}{\hookrightarrow} Y, \qquad
i_{!}(\mathbb{Q}_{\ell,X_{0}}) \to Ri_{*}(\mathbb{Q}_{\ell,X_{0}}) .
\]

(1) $X$ \emph{régulier}, $F = \mathbb{Q}_{\ell,X}$, $f$ immersion
($n = \dim \overline{f(X)}$).
On peut supposer $X \to Y$ imm.\ \emph{ouverte}, $Y'$ régulier,
$Y' - g(X) = h^{-1}(Y - X) = Z$ diviseur à croisements normaux,
$Z = \sum Z_{i}$ :

\begin{tikzcd}
X \arrow[r, "g"] \arrow[dr, "f"'] & Y' \arrow[d, "h"] \\
& Y
\end{tikzcd}

\[
R^{*}f_{*}(\mathbb{Q}_{\ell,X}) \Longleftarrow E_{2}^{pq} =
R^{p}h_{*}\bigl(R^{q}g_{*}(\mathbb{Q}_{\ell})\bigr),
\]
\[
R^{q}g_{*}(\mathbb{Q}_{\ell}) = \bigwedge^{q} \sum \mathbb{Q}_{\ell,Z_{i}} =
\sum_{i_{1} < \cdots < i_{q}} \mathbb{Q}_{\ell,Z_{i_{1}} \cap \cdots \cap Z_{i_{q}}}(-q)
\]
pur de poids $2q$ ;

$E_{2}^{p0} = R^{p}h_{*}(\mathbb{Q}_{\ell,Y'})$ est eff\supplied{t} admissible de
poids $\leqslant p = i$, niveau
$\leqslant \operatorname{Inf}(p, 2n-2-p)$ [avec \uncertain{égalité} si
$p \geqslant n-1$] ;

$E_{2}^{pq} = \sum_{i_{1} < \cdots < i_{q}}
R^{p}h_{Z_{i_{1}} \cap \cdots \cap Z_{i_{q}}\,*}(\mathbb{Q}_{\ell})(-q)$ est
eff.\ admissible de poids $\leqslant p + 2q = i + q$
\struck{$\leqslant 2n$}, de niveau $\leqslant \operatorname{Inf}(p, 2n-2q-p)$.
\note{les indices des faisceaux sont écrits $\mathbb{Q}_{\ell,X_{0}}$ puis
récrits ; celui de $E_{2}^{p0}$ est surchargé, $Y'$ \uncertain{probable}}

En tous cas, poids sont $\leqslant \operatorname{Inf}(2i, 2n-2 \ldots$
\note{l'encadré de « $\leqslant$ » est surchargé}, niveau
$\leqslant \operatorname{Inf}(i, 2n-2-i)$.

\marginal{(5.e) \quad $i+q$ : \quad $p \leqslant 2(n-q)$, encadrés
$p + 2q \leqslant 2n$ et $q \leqslant n, 2n-i$ ; puis $p \leqslant n-1$,
$q \leqslant n$}

\note{en bas de page, des calculs d'indices : $p$, $n-q$ ; $p+q \geqslant n$ ;
$2(p-n-q) = 2(i-n)$ ; $2(n-q) - p = 2n - i - q$ ; et, à droite, un
diagramme des $(p, q)$ admissibles, points marqués d'une croix dans le
quadrant, axes gradués $1, 2, \ldots, n-1, n$ et $n-3, n-2, n-1$, que nous
ne redessinons pas}

\page{117}
\note{feuillet de tableaux de poids et de niveaux, suite des calculs de la
page 115 ; il n'y a pas de phrases. Nous donnons les tableaux
dans l'ordre de la page, en omettant les colonnes intermédiaires qu'il
laisse en pointillé}

$i \leqslant n-1$ :
\[
\begin{array}{c|ccc}
\text{poids} & 0 & 1 \;\cdots & 2i \\
\hline
\text{niveau} & 0 & 1 \;\cdots & 2i
\end{array}
\]
\note{premier essai, séparé du reste par un trait ; le niveau sous $2i$ est
« $2i$ », peut-être corrigé plus bas en 0}

\[
\begin{array}{l|c|c|c|c|c|c}
E_{2}^{i,0} & \text{poids} & 0 & 1 & \cdots & i & \\
 & \text{niveau} \leqslant & 0 & 1 & \cdots & i & \\
\hline
E_{2}^{i-1,1} & \text{poids} & 2 & 3 & \cdots & i & i+1 \\
 & \text{niveau} & 0 & 1 & \cdots & i-2 & i-1 \\
\hline
E_{2}^{i-2,2} & \text{poids} & 4 & 5 & \cdots & i \;\; i+1 & i+2 \\
 & \text{niveau} & 0 & 1 & \cdots & i-4 \;\; i-3 & i-2 \\
\hline
E_{2}^{0,i} & \text{poids} & 2i & & & & \\
 & \text{niveau} & 0 & & & &
\end{array}
\]
\note{en marge de ces lignes, des conditions biffées :
\struck{$p+q = i$}, \struck{$p \leqslant 2(n-1-q)$}, \struck{$p+q \leqslant
n-1$} ; la dernière entrée de la ligne $E_{2}^{i-1,1}$ est surchargée}

$i \leqslant n-1$ (encadré) :
\[
\begin{array}{c|c|c|c|c|c|c|c|c}
\text{poids} & 0 & 1 & \cdots & i-1 & i & i+1 & \cdots & 2i \\
\hline
\text{niveau} & 0 & 1 & \cdots & i-1 & i & i-1 & \cdots & 0
\end{array}
\]

$i \geqslant n$ :
\[
\begin{array}{l|c|c|c|c|c}
E_{2}^{i,0} & 2(i-n+1) & 2(i-n+1)+1 & \cdots & & i \\
 & 0 & 1 & 2 \;\cdots & & 2(n-1)-i \\
\hline
E_{2}^{i-1,1} & 2(i-n)+2 & 2(i-n)+3 & \cdots & i & i+1 \\
 & 0 & 1 & \cdots & 2n-i-2 & 2n-i-1 \\
\hline
E_{2}^{i-2,2} & 2(i-n)+4 & 2(i-n)+5 & \cdots & i \;\; i+1 & i+2 \\
 & 0 & 1 & \cdots & 2n-i-2 \;\; 2n-i-3 & 2n-i-2 \\
\hline
E_{2}^{i-3,3} & 2(i-n)+6 & 2(i-n)+7 & \cdots & i \;\; i+1 \;\; i+2 & i+3 \\
 & 0 & 1 & \cdots & 2n-i-6 \;\; \ldots & 2n-i-3 \\
\hline
E_{2}^{0,i} & 2i & & & & \\
 & 0 & & & &
\end{array}
\]
\note{les niveaux des dernières colonnes sont écrits vite et en partie
surchargés ; nous les donnons tels que nous les lisons, sans les
vérifier. Des traits obliques relient la dernière colonne de chaque ligne à
la suivante}

\[
q = 2n - i, \quad p = 2(i-n) : \qquad
E_{2}^{2(i-n),\,2n-i} = R^{2(i-n)}f_{Z\,*}(\mathbb{Q}_{\ell}\ldots, \qquad
2(i-n) + 2(2n-i) = 2n .
\]

\note{à gauche, deux diagrammes des couples $(p, q)$ : des croix dans le
quadrant, axes gradués $0, 1, 2, 3, \ldots, n-2, n-1, n$ en ordonnée et
$\ldots, 2n-6, 2n-5, 2n-4, 2n-3, 2n-2$ en abscisse, séparés par un trait
vertical ; nous ne les redessinons pas}

\page{119}
\note{feuillet de calculs, qui reprend les tableaux de la page
117 ; il n'y a guère de phrases}

$i \geqslant n$ (encadré) :
\[
\begin{array}{c|c|c|c|c|c|c|c|c}
\text{poids} & 2(i-n)+2 & 2(i-n)+3 & 2(i-n)+4 & \cdots & i+1 & i+2 & \cdots & 2n \\
\hline
\text{niveau} & 0 & 1 & 2 & \cdots & 2n-i-1 & 2n-i-2 & \cdots & 0
\end{array}
\]
\note{une colonne « $i$ », de niveau $2n-i-2$, est barrée d'une croix entre
$2(i-n)+4$ et $i+1$ ; sous le tableau, deux flèches, l'une montante vers le
milieu, l'autre descendante : les niveaux croissent puis décroissent}

Poids $2n$ intervient avec
\[
E_{2}^{2(i-n),\,2n-i} = \sum R^{2(i-n)}f_{Z_{\alpha_{1}} \cap Z_{\alpha_{2}}
\cap \cdots \cap Z_{\alpha_{2n-i}}\,*}(\mathbb{Q}_{\ell})\bigl(-(2n-i)\bigr),
\]
l'intersection étant de dim.\ $i-n$ ; $2n-i \leqslant n \leqslant i$ ;
$n - (2n-i) = i-n$ ; $2(i-n) + 2(2n-i) = 2n$.

\[
\begin{array}{c|c|c|c|c|c|c|c|c|l}
i & 0 & 1 & \cdots & n-1 & n & n+1 & \cdots & 2n-1 & \\
\hline
\text{niveau} & 0 & 1 & \cdots & n-1 & n-1 & n-2 & \cdots & 0 & \leqslant n-1 \\
\hline
\text{poids} \leqslant & 0 & 2 & \cdots & 2n-2 & 2n & 2n & 2n & 2n & \leqslant 2n
\end{array}
\]
\note{tableau encadré ; en tête de la dernière colonne, « $\leqslant 2n-1$ »,
et un « $\leqslant 2n$ » biffé}

\[
\boxed{\begin{array}{l}
\text{poids} \leqslant 2\operatorname{Inf}(i, n) \\
\text{niveau} \leqslant \operatorname{Inf}(i, n-1)
\end{array}}
\]

Si $F$ \uncertain{str\supplied{t}} spécial de poids $j$ :
\[
\text{poids} \leqslant 2\operatorname{Inf}(i, n) + j \;? \qquad
\text{niveau} \leqslant \operatorname{Inf}(i, n-1) + j .
\]

(2°) \struck{$f$ \ill{} $F = \mathbb{Q}_{\ell,X}$} Cas général
\uncertain{(relatif, en gardant \ill{})} :

\begin{tikzcd}
X \arrow[d, "f"] \\
Y
\end{tikzcd}

$y \in Y$, et
\[
n = d_{y}(f) = \operatorname*{Sup}_{x \in X} \bigl[\deg \operatorname{tr}
k(x) : k(f(x)) + \dim \overline{f(x)}_{y}\bigr] ,
\]
$x$ tel que $f(x) \in \uncertain{\operatorname{Spec} \mathcal{O}_{Y,y}}$ ;
\emph{au voisinage de $y$} :
\[
\boxed{\begin{array}{l}
\text{poids} \leqslant 2\operatorname{Inf}(i, n) + j \\
\text{niveau} \leqslant \operatorname{Inf}(i, n) + j
\end{array}}\ ?
\]
\marginal{peut-on remplacer $n$ par \uncertain{la} \struck{dim des fibres}
… $\leqslant n-1$ ?}
\note{la formule de $d_{y}(f)$ est écrite vite et en partie surchargée :
lecture douteuse dans le détail}

\page{120}
\note{feuillet dactylographié, en français, isolé : il commence et s'arrête au
milieu d'une phrase, et ne continue ni ne précède rien dans ce lot. Il
présente des travaux de Deligne. Les soulignements, et les ajouts « resté »,
« sur X » et « -Puppe », sont à la main ; les biffures à la machine sont
rendues par \texttt{struck} quand elles se lisent}

une sorte de synthèse géométrico-arithmétique des nombreuses théories
cohomologiques dont on dispose \struck{\ill{}} à présent pour les variétés
algébriques, \struck{de toutes dimensions,} et des représentations linéaires
de groupes de Galois arithmétiques auxquelles ces théories donnent naissance.

5) \emph{Solution du problème d'introduction d'une $\lambda$-structure}
(correspondant à la notion de puissance extérieure) \emph{dans l'anneau}
$K^{\bullet}(X)$ (dit « de Grothendieck ») d'un espace topologique localement
annelé (par exemple une variété analytique, ou différentiable, ou un schéma),
construit à l'aide de la catégorie des complexes « parfaits » sur $X$. Ce
problème était signalé comme le plus important problème \add{resté} en
suspens dans le Séminaire de Géométrie Algébrique de l'IHES 1966/67, sur le
théorème de Riemann-Roch. La nécessité de travailler avec le $K^{\bullet}(X)$
défini en termes de \emph{complexes} parfaits, plutôt que de \emph{fibrés}
vectoriels seulement, est imposé par des \struck{\ill{}} impératifs
techniques irrécusables, qui finiront certainement par s'imposer aussi en
dehors de la géométrie algébrique. La solution de Deligne, comme prévu, passe
par l'intermédiaire d'une construction de \struck{\ill{}} foncteurs puissances
extérieures dans la catégorie dérivée $D(X)$ (catégorie des complexes de
modules \add{sur $X$} --- en fait, il faut se borner ici aux complexes de
\emph{chaînes}), en utilisant la théorie de Dold\supplied{-Puppe}. La difficulté
sérieuse consistait à prouver que ces opérations dans $K^{\bullet}(X)$
satisfont aux propriétés algébriques habituelles, permettant d'exprimer par
exemple $\lambda^{i}(xy)$ comme polynôme universel à coéfficients entiers en
les $\lambda^{j}(x), \lambda^{k}(y)$. Deligne procède par une voie très
ingénieuse, en déterminant d'abord la structure du $\lambda$-anneau de toutes
les « opérations tensorielles universelles » qu'on peut effectuer sur des
fibrés vectoriels. Il n'a écrit sa théorie que lorsqu'on travaille au dessus
d'un corps de base \struck{\ill{}} (de caractéristique arbitraire).
\struck{Le cas de la caractéristique}
\note{la dernière ligne est biffée à la main ; le feuillet s'arrête là}

\end{document}
